\section{Del pipeline vertical al sequence learning general}

Cuando la línea histórica llega al NLP, el cambio más evidente no es el aumento de parámetros de los modelos, sino la reducción de las interfaces del sistema. La traducción automática de 2009 requería word alignment, phrase extraction, reordenamiento, language-model rescoring y una gran cantidad de características independientes; en 2013, la comunidad de investigación también se organizaba por áreas verticales como morphology, dialogue, discourse, sentiment y translation. La sequence formulation general fue convirtiendo poco a poco estas tareas en un problema compartido de aprendizaje de representaciones.

\subsection{2009: la red neuronal es solo un rescorer dentro de un sistema complejo}

La statistical machine translation (SMT, traducción automática estadística) temprana armaba un sistema completo combinando varios módulos interpretables. En esa época, el neural probabilistic language model solía colocarse al final del pipeline para volver a puntuar las traducciones candidatas, no para decidir la salida de extremo a extremo. Una ingeniería así podía funcionar, pero cada interfaz tenía que definir representación, búsqueda, características y propagación de errores, y el conocimiento quedaba encerrado en componentes especializados.

\lecturefigure{slide-02-sequence-models-2009.jpg}{Pipeline de traducción automática de 2009 y neural probabilistic language model}{Diapositiva V002 recuperada del video oficial; 00:05:09}

\begin{knowledgebox}{Lectura de la figura: dónde está realmente la red neuronal}
La imagen de la izquierda incluye varios módulos, como alignment, phrase table, decoder y rescoring; el neural language model de la derecha solo aporta una de esas puntuaciones. Primero cuente las fronteras del sistema y después observe si el gradiente puede propagarse de extremo a extremo a través de ellas. La figura respalda que el neural component ya tenía valor en ese momento, pero no controlaba la tarea completa, y también explica por qué una consolidation basada en datos traería enormes beneficios de ingeniería.
\end{knowledgebox}

\teachervoice{Vaswani recuerda que el sistema de MT de cuando hacía el doctorado era aún más complejo que el de la figura, y pide a todo el grupo que encuentre «dónde está la red neuronal». La respuesta es que el Continuous State Language Model se usaba principalmente para rescoring; vista desde hoy, la insistencia del nombre en continuous indica precisamente que el mundo dominante de entonces seguía siendo el del pipeline discreto.}

\subsection{2013: las tareas y las comunidades de investigación se verticalizan}

Cuando los modelos no pueden compartir suficiente representación y estructura de cómputo, cada problema necesita su propio conocimiento del dominio, sus propias características y su propia evaluation culture. Un \term{verticalized system (sistema verticalizado)} es una pila tecnológica construida alrededor de una tarea, con datos, módulos y métricas especializados; permite profundizar en el problema local, pero dificulta la transfer entre tareas y la infraestructura común.

\lecturefigure{slide-03-nlp-2013-verticalized.jpg}{Las líneas de investigación de NLP en 2013 muestran una fuerte verticalización}{Diapositiva V003 recuperada del video oficial; 00:07:03}

\begin{knowledgebox}{Lectura de la figura: el índice de un congreso también es un organigrama técnico}
La página enumera varios research track relativamente independientes. Primero observe en cuántas tareas se divide un mismo fenómeno lingüístico y luego pregúntese si cada tarea usa feature, modelo y datos propios. Esto no demuestra que la investigación especializada carezca de valor; más bien sugiere que la consolidation necesita una interfaz de cómputo general capaz de manejar en conjunto variable-length input, context y conditional generation.
\end{knowledgebox}

\subsection{Distributed Representation, Seq2Seq y Attention}

La \term{distributed representation (representación distribuida)} codifica un concepto con varias dimensiones de un vector denso en conjunto, en lugar de mantener una característica simbólica independiente para cada palabra o regla; la \term{contextual representation (representación contextualizada)} va más allá y hace que el vector de un mismo token cambie según el contexto. Seq2Seq unifica tareas como translation, summarization, dialogue y question answering bajo la forma «secuencia de entrada a secuencia de salida», y attention permite que el decoder lea los source states según su contenido.

\lecturefigure{slide-04-general-approaches-2016.jpg}{El paradigma general: de los vectores de palabras a Seq2Seq y Attention}{Diapositiva V004 recuperada del video oficial; 00:09:21}

\begin{knowledgebox}{Términos clave: tres ascensos de abstracción}
\begin{tabularx}{\textwidth}{@{}L{0.20\textwidth}X X@{}}
\toprule
Abstracción & Problema que resuelve & Cuello de botella que persiste \\
\midrule
Distributed word vectors & Comparte la estructura estadística de las palabras y la similitud semántica & Una misma palabra sigue siendo bastante fija en contextos distintos \\
Contextual sequence states & Hace que la representación dependa del contexto anterior y posterior & La recurrent computation sigue siendo secuencial en el tiempo \\
Seq2Seq + attention & Unifica la generación condicional y lee la source según su contenido & La columna vertebral encoder/decoder suele seguir formada por RNN \\
\bottomrule
\end{tabularx}
\end{knowledgebox}

\teachervoice{El ponente describe esta evolución como «general paradigms started coming up». Lo importante no es lo llamativa que resulte el álgebra vectorial de Word2Vec, sino que los investigadores empezaron a creer que, si se lograba aprender de forma estable una variable-length representation, muchas tareas lingüísticas antes independientes podrían compartir una misma forma de modelado.}

\subsection{La variable-length representation es el núcleo común}

Un \term{sequence model (modelo de secuencias)} procesa entradas o salidas ordenadas y de longitud variable. Debe convertir token sequences de longitudes distintas en estados utilizables para la predicción, conservando a la vez el contenido, el orden y las dependencias de largo alcance. La figura siguiente usa un recurrent encoder para representar la ruta clásica: el estado se actualiza paso a paso y el decoder y las tareas posteriores usan la representación final o la de cada posición.

\lecturefigure{slide-05-variable-length-representations.jpg}{Aprender variable-length representations es la base común de las tareas Seq2Seq}{Diapositiva V005 recuperada del video oficial; 00:10:33}

Una recurrent update mínima es:
\[
h_t=f_\theta(h_{t-1},x_t),\qquad y_t\sim p_\theta(\cdot\mid h_t).
\]
Aquí $x_t$ es el $t$-ésimo token de entrada, $h_t$ es el hidden state que resume el prefijo $x_{1:t}$, $f_\theta$ es la recurrent transition compartida y $y_t$ es la predicción. La recursión permite que el modelo exprese el orden de forma natural, pero obliga a que $h_t$ espere a $h_{t-1}$, por lo que la profundidad de paralelización crece con la longitud de la secuencia.

\begin{importantbox}{El primer paso de la consolidation es unificar la formulación del problema}
Antes del Transformer, Seq2Seq ya había logrado una abstracción importante: translation, summarization y dialogue podían escribirse como conditional sequence modeling. La aportación del Transformer no fue inventar la unificación desde cero, sino liberar esa interfaz unificada del sequential backbone.
\end{importantbox}

\subsection{Resumen de la sección}

El pipeline de MT de 2009 y las líneas de investigación verticales de 2013 muestran que la specialization llegó a penetrar cada interfaz del sistema. La distributed/contextual representation, Seq2Seq y attention unificaron primero la forma de las tareas; la siguiente pregunta no es si las RNN tienen capacidad expresiva, sino si esa capacidad expresiva puede entrenarse con eficiencia a gran escala de datos y de hardware.
